\subsection{Representation}\label{representation}

Let $\boldsymbol{x}_t = [x_{1,t}, x_{2,t}, \dots, x_{N,t}]'$, for $t = 1, \dots, T$, denote a stationary $N$-dimensional vector process standardised to have mean zero and unit variance. We assume that $\boldsymbol{x}_t$ admits the following factor model representation:
\begin{equation}
\boldsymbol{x}_t = \boldsymbol{\Lambda} \boldsymbol{f}_t + \boldsymbol{e}_t
\end{equation}
where $\boldsymbol{f}_t$ is an $r \times 1$ vector of unobserved common factors, $\boldsymbol{\Lambda}$ is an $N \times r$ matrix of factor loadings, and $\boldsymbol{e}_t = [e_{1,t}, \dots, e_{N,t}]'$ denotes the idiosyncratic component. The term $\boldsymbol{\chi}_t = \boldsymbol{\Lambda} \boldsymbol{f}_t$ is referred to as the common component.

The common factors follow a stationary vector autoregressive process of order $p$:
\begin{equation}
\boldsymbol{f}_t = \sum_{j=1}^{p} \boldsymbol{\Phi}_j \boldsymbol{f}_{t-j} + \boldsymbol{\eta}_t
\end{equation}
where $\boldsymbol{\eta}_t$ is a vector of mean-zero innovations with covariance matrix $\boldsymbol{\Sigma}_{\eta}$. The idiosyncratic components $\boldsymbol{e}_t$ are allowed to exhibit weak cross-sectional and serial dependence, giving rise to an approximate factor model. It is assumed that $\boldsymbol{e}_t$ is uncorrelated with $\boldsymbol{\eta}_t$ at all leads and lags.

\subsection{Estimation}

We consider three approaches to estimate the dynamic factor model: principal components analysis (PCA), quasi-maximum likelihood (QML) estimation, and a two-step estimator. PCA provides a simple and robust benchmark by estimating the factor space from the covariance structure of the data without imposing an explicit dynamic model. QML exploits the state-space representation and incorporates factor dynamics, improving efficiency under stronger assumptions. The two-step estimator combines these approaches by retaining the tractability of PCA while incorporating dynamic structure through state-space methods.

\paragraph{Principal components.}

PCA is a nonparametric approach, as it does not require specifying a model for the factor dynamics or imposing distributional assumptions on the idiosyncratic components. Under the static representation, the factor loadings $\boldsymbol{\Lambda}$ and the factors $\boldsymbol{f}_t$ can be estimated by minimising:
\begin{equation}
SSR(\boldsymbol{\Lambda}, \boldsymbol{f}_t) = \frac{1}{NT} \sum_{t=1}^{T} (\boldsymbol{x}_t - \boldsymbol{\Lambda} \boldsymbol{f}_t)' (\boldsymbol{x}_t - \boldsymbol{\Lambda} \boldsymbol{f}_t)
\end{equation}

Due to the indeterminacy of the factors and loadings, identification requires normalisation.\footnote{This reflects the fact that both the factors and loadings are unobserved and only their product is identified. To illustrate the identification issue, consider any invertible $r \times r$ matrix $\boldsymbol{Q}$. Then, the common component $\boldsymbol{\Lambda} \boldsymbol{f}_t$ is observationally equivalent to $\boldsymbol{\Lambda}^* \boldsymbol{f}_t^* = (\boldsymbol{\Lambda} \boldsymbol{Q}^{-1})(\boldsymbol{Q} \boldsymbol{f}_t)$, where $\boldsymbol{\Lambda}^* = \boldsymbol{\Lambda} \boldsymbol{Q}^{-1}$ and $\boldsymbol{f}_t^* = \boldsymbol{Q} \boldsymbol{f}_t$.} A convenient normalisation imposes $N^{-1}\boldsymbol{\Lambda}'\boldsymbol{\Lambda} = \boldsymbol{I}_r$ and assumes that $\mathbb{E}(\boldsymbol{f}_t \boldsymbol{f}_t')$ is diagonal.

Under these conditions, the PCA estimators are obtained from the eigenvalue decomposition of the sample covariance matrix:
\begin{equation}
\hat{\boldsymbol{\Sigma}}_x = \frac{1}{T} \sum_{t=1}^{T} \boldsymbol{x}_t \boldsymbol{x}_t'
\end{equation}
and correspond to the eigenvectors associated with the largest $r$ eigenvalues. The estimated factors are then given by:
\begin{equation}
\hat{\boldsymbol{f}}_t = N^{-1} \hat{\boldsymbol{\Lambda}}' \boldsymbol{x}_t
\end{equation}

The consistency of the PCA estimator was established by \citet{bai2002determining} and \citet{stock2002forecasting}, while \citet{bai2003inferential} provides asymptotic distribution theory. After estimating the factors, their dynamics can be characterised by fitting a VAR model. Based on standard information criteria, we select $p = 4$.

\paragraph{Quasi-maximum likelihood (QML) estimation.}

PCA does not explicitly account for factor dynamics, which may result in efficiency losses in estimation. QML addresses this by exploiting the state-space representation of the model, combining cross-sectional and time-series information.

The model can be written in state-space form with measurement equation:
\begin{equation}
\boldsymbol{x}_t = \boldsymbol{\Lambda} \boldsymbol{f}_t + \boldsymbol{e}_t
\end{equation}
and transition equation:
\begin{equation}
\boldsymbol{F}_t = \boldsymbol{A} \boldsymbol{F}_{t-1} + \boldsymbol{B} \boldsymbol{\eta}_t
\end{equation}
where the state vector is defined as:
\begin{equation}
\boldsymbol{F}_t =
\begin{bmatrix}
\boldsymbol{f}_t \\
\boldsymbol{f}_{t-1} \\
\vdots \\
\boldsymbol{f}_{t-p+1}
\end{bmatrix}
\end{equation}
and the system matrices $\boldsymbol{A}$ and $\boldsymbol{B}$ are given by:
\begin{equation}
\boldsymbol{A} =
\begin{bmatrix}
\boldsymbol{\Phi}_1 & \cdots & \boldsymbol{\Phi}_p \\
\boldsymbol{I}_r & \cdots & \boldsymbol{0} \\
\vdots & \ddots & \vdots \\
\boldsymbol{0} & \cdots & \boldsymbol{I}_r
\end{bmatrix},
\qquad
\boldsymbol{B} =
\begin{bmatrix}
\boldsymbol{I}_r \\
\boldsymbol{0} \\
\vdots \\
\boldsymbol{0}
\end{bmatrix}
\end{equation}

In large cross-sectional settings, full maximum likelihood estimation becomes computationally demanding. \citet{doz2012quasi} propose a QML approach based on a simplified state-space representation, which assumes that the factors follow a Gaussian distribution and that the idiosyncratic components are mutually orthogonal. Estimation proceeds via the Expectation-Maximisation (EM) algorithm \citep{watson1983alternative}, which alternates between a Kalman filter-based expectation step and a maximisation step based on multivariate regressions.

Importantly, the estimator is consistent as both $N$ and $T$ increase, even when these distributional assumptions are relaxed, including in the presence of non-Gaussian factors and weak cross-sectional dependence in the idiosyncratic components. This robustness underpins its interpretation as a \textit{quasi}-maximum likelihood estimator.

\paragraph{Two-step estimator.}

A computationally attractive alternative is the two-step estimator proposed by \citet{doz2011two}, which combines PCA and state-space methods.

In the first step, preliminary estimates $\hat{\boldsymbol{\Lambda}}$ and $\hat{\boldsymbol{f}}_t$ are obtained via PCA, and a VAR model is fitted to the estimated factors to obtain $\hat{\boldsymbol{\Phi}}_j$. In the second step, the model is cast in state-space form assuming $\text{cov}(\boldsymbol{\eta}_t) = \boldsymbol{I}_r$ and a diagonal covariance matrix for $\boldsymbol{e}_t$, and a Kalman smoother is applied to obtain refined factor estimates:
\begin{equation}
\hat{\boldsymbol{f}}_{t|T} = \mathbb{E}(\boldsymbol{f}_t \mid \boldsymbol{x}_1, \dots, \boldsymbol{x}_T)
\end{equation}

Although the idiosyncratic components are treated as mutually orthogonal in the state-space representation, \citet{doz2011two} show that the resulting factor estimates remain consistent, as the bias induced by this misspecification vanishes asymptotically.

\paragraph{Asymptotic distribution and confidence bands.}

To characterise the uncertainty in the estimated financial cycle, we draw on the inferential theory of \citet{bai2003inferential}. Under standard regularity conditions on the factor model and as both $N$ and $T$ tend to infinity, the PCA estimator of the common factor satisfies the following limiting distribution:
\begin{equation}
    \sqrt{N}\left(\hat{\boldsymbol{f}}_t - \boldsymbol{H}\boldsymbol{f}_t^0\right) \xrightarrow{d} \mathcal{N}\!\left(0,\; V_{NT}^{-2}\,\hat{\Gamma}_t\right),
    \label{eq:factor_clt}
\end{equation}
where $\boldsymbol{H}$ is an asymptotically invertible rotation matrix, $V_{NT}$ denotes the first eigenvalue of $(NT)^{-1}\boldsymbol{X}\boldsymbol{X}'$, and $\Gamma_t$ is the asymptotic variance of the cross-sectional average of the weighted idiosyncratic components. A consistent estimator of the asymptotic variance is given by \citep{bai2003inferential,zhang2019factor}:
\begin{equation}
    \widehat{\mathrm{Avar}}(\hat{\boldsymbol{f}}_t) = V_{NT}^{-2}\,\hat{\Gamma}_t,
    \label{eq:avar}
\end{equation}
where $\hat{\Gamma}_t$ is estimated as:
\begin{equation}
    \hat{\Gamma}_t = \frac{1}{N}\sum_{i=1}^{N} \hat{e}_{it}^2\,\hat{\lambda}_i^2,
    \label{eq:gamma_est}
\end{equation}
with $\hat{e}_{it} = x_{it} - \hat{\lambda}_i\hat{f}_t$ denoting the estimated idiosyncratic component and $\hat{\lambda}_i$ the PCA loading of variable $i$. This estimator is consistent under the assumption that the idiosyncratic errors are cross-sectionally uncorrelated but may exhibit heteroskedasticity across both variables and time.\footnote{Cross-sectionally and heteroskedastically autocorrelation consistent (CS-HAC) estimators of $\Gamma_t$, which allow for cross-sectional dependence in $e_{it}$, are algebraically degenerate in the PCA setting: by the first-order conditions of the PCA objective, the estimated loadings $\hat{\boldsymbol{\lambda}}$ are orthogonal to the residuals $\hat{\boldsymbol{e}}_t$, so the full CS-HAC sum collapses to zero identically.}

An $(1-\alpha)$ confidence interval for $\boldsymbol{H}\boldsymbol{f}_t^0$ at each $t$ is then:
\begin{equation}
    \hat{f}_t \;\pm\; z_{\alpha/2}\,\sqrt{\widehat{\mathrm{Avar}}(\hat{f}_t)} \;=\; \hat{f}_t \;\pm\; z_{\alpha/2}\,\frac{N\sqrt{\hat{\Gamma}_t}}{\lambda_1},
    \label{eq:ci}
\end{equation}
where $\lambda_1 = N\,V_{NT}$ is the first eigenvalue of the sample covariance matrix of the data and $z_{\alpha/2}$ is the corresponding normal quantile.\footnote{The factor estimator used in practice is $\hat{f}_t = \hat{\boldsymbol{\lambda}}'\boldsymbol{x}_t$, where $\hat{\boldsymbol{\lambda}}$ is the unit-norm PCA eigenvector. This differs from the PC2-normalised factor in \citet{bai2003inferential} by a factor of $\sqrt{N}$, so its asymptotic standard error is $\sqrt{N}$ times larger: $\mathrm{SE}(\hat{f}_t) = \sqrt{\widehat{\mathrm{Avar}}(\hat{f}_t)}$ without the additional $1/\sqrt{N}$ correction.} We construct 95\% confidence bands ($\alpha = 0.05$, $z_{0.025} \approx 1.96$). The bands are centred on the financial cycle indicator at each point in time. Periods in which the indicator lies within the confidence band are classified as \textit{neutral}: the factor is not significantly different from its long-run average of zero at the 5\% significance level. Since $\hat{\Gamma}_t$ varies over time through the squared residuals $\hat{e}_{it}^2$, the bands widen during periods of large idiosyncratic variation and tighten when the common factor explains a greater share of cross-sectional variation. The mean of the upper and lower bounds provides a pair of constant dashed-line thresholds used for the regime classification throughout the paper.
